\documentclass[11pt]{article}
\usepackage[margin=1in]{geometry}
\usepackage{amsmath,amssymb,amsthm,booktabs}
\newtheorem{theorem}{Theorem}
\newtheorem{proposition}{Proposition}
\title{Dependency-Aware Local Typical Sets and a Two-Center $SU(2)$ Atlas}
\author{}
\date{}
\begin{document}
\maketitle

\begin{abstract}
We replace a global high-dimensional Gaussian ball by local link, plaquette
input, and curvature constraints.  The resulting coverage estimate requires
no independence.  For the finite $2^4$ Coulomb Gaussian, the required
four-link radius is approximately $0.00602316$, reducing the apparent chart
gap by a factor of four.  We also record an exact two-center atlas for all of
$SU(2)$.
\end{abstract}

\section{Local Gaussian coverage}
Let $A_j$ be jointly Gaussian blocks in $\mathbb R^r$ and suppose their
marginal covariance operators satisfy $\operatorname{Cov}(A_j)\preceq
\lambda I$.  No assumption is made on correlations between different blocks.

\begin{theorem}[Local union bound]
For $N$ blocks and $0<\delta<1$,
\[
 \mathbb P\left\{\max_{1\le j\le N}\|A_j\|>
 \sqrt\lambda\left(\sqrt r+\sqrt{2\log(N/\delta)}\right)\right\}
 \le\delta.
\]
\end{theorem}
\begin{proof}
Write each marginal as $\Sigma_j^{1/2}Z$ with $Z$ standard Gaussian.  Since
$\|\Sigma_j^{1/2}\|\le\sqrt\lambda$, Gaussian norm concentration gives
$\mathbb P(\|A_j\|>\sqrt\lambda(\sqrt r+\sqrt{2x}))\le e^{-x}$.
Take $x=\log(N/\delta)$ and apply the union bound.  Dependence never enters.
\end{proof}

On the scalar Coulomb space $\ker d_0^T$ of the $2^4$ block, the Hessian
$d_1^Td_1$ has physical dimension $136$ and spectrum in $[2,12]$.  Direct
finite spectral calculation gives the leverage coefficients shown below.

\begin{center}
\begin{tabular}{lrr}
\toprule
Local object & covariance coefficient & radius at $t=3\cdot10^{-6}$\\
\midrule
one link ($r=3$) & $0.128788574$ & $0.004025831$\\
four-link plaquette input ($r=12$) & $0.179460637$ & $0.006023158$\\
linear curvature ($r=3$) & $0.685582011$ & $0.009288521$\\
\bottomrule
\end{tabular}
\end{center}
Here $N=216$ and $\delta=0.00279270808$.  By comparison, the global
$408$-dimensional ball of radius $0.003$ had log mass $-803.25266$; it was a
valid but structurally mismatched certificate.

\section{Two-center group atlas}
Identify $SU(2)$ with unit quaternions.  Choose the nearer center $z\in\{1,-1\}$
according to the sign of the scalar quaternion component and write
$zU=\exp X$.  Then $\|X\|\le\pi/2$.  Thus the two center charts cover $SU(2)$
with multiplicity at most two.  In exponential coordinates the normalized
Haar density is
\[
 J(r)=\frac{1}{2\pi^2}\left(\frac{\sin r}{r}\right)^2,
\]
so its ratio on either radius-$\pi/2$ chart is at most $(\pi/2)^2$.
Relative tree gauge has constant gauge-volume Jacobian.

\section{Scope}
The local-union theorem and group atlas do not yet compare normalized
nonlinear and Gaussian conditional measures.  A charged estimate for the
nonprincipal center and orbit channels is also absent.  The next question is
whether nonlinear quotient convexity reaches the radius $0.006023158$.
\end{document}

